% Options for packages loaded elsewhere
\PassOptionsToPackage{unicode}{hyperref}
\PassOptionsToPackage{hyphens}{url}
%
\documentclass[
]{article}
\usepackage{amsmath,amssymb}
\usepackage{lmodern}
\usepackage{iftex}
\ifPDFTeX
  \usepackage[T1]{fontenc}
  \usepackage[utf8]{inputenc}
  \usepackage{textcomp} % provide euro and other symbols
\else % if luatex or xetex
  \usepackage{unicode-math}
  \defaultfontfeatures{Scale=MatchLowercase}
  \defaultfontfeatures[\rmfamily]{Ligatures=TeX,Scale=1}
\fi
% Use upquote if available, for straight quotes in verbatim environments
\IfFileExists{upquote.sty}{\usepackage{upquote}}{}
\IfFileExists{microtype.sty}{% use microtype if available
  \usepackage[]{microtype}
  \UseMicrotypeSet[protrusion]{basicmath} % disable protrusion for tt fonts
}{}
\makeatletter
\@ifundefined{KOMAClassName}{% if non-KOMA class
  \IfFileExists{parskip.sty}{%
    \usepackage{parskip}
  }{% else
    \setlength{\parindent}{0pt}
    \setlength{\parskip}{6pt plus 2pt minus 1pt}}
}{% if KOMA class
  \KOMAoptions{parskip=half}}
\makeatother
\usepackage{xcolor}
\usepackage{longtable,booktabs,array}
\usepackage{multirow}
\usepackage{calc} % for calculating minipage widths
% Correct order of tables after \paragraph or \subparagraph
\usepackage{etoolbox}
\makeatletter
\patchcmd\longtable{\par}{\if@noskipsec\mbox{}\fi\par}{}{}
\makeatother
% Allow footnotes in longtable head/foot
\IfFileExists{footnotehyper.sty}{\usepackage{footnotehyper}}{\usepackage{footnote}}
\makesavenoteenv{longtable}
\setlength{\emergencystretch}{3em} % prevent overfull lines
\providecommand{\tightlist}{%
  \setlength{\itemsep}{0pt}\setlength{\parskip}{0pt}}
\setcounter{secnumdepth}{-\maxdimen} % remove section numbering
\ifLuaTeX
  \usepackage{selnolig}  % disable illegal ligatures
\fi
\IfFileExists{bookmark.sty}{\usepackage{bookmark}}{\usepackage{hyperref}}
\IfFileExists{xurl.sty}{\usepackage{xurl}}{} % add URL line breaks if available
\urlstyle{same} % disable monospaced font for URLs
\hypersetup{
  hidelinks,
  pdfcreator={LaTeX via pandoc}}

\author{}
\date{}

\begin{document}

\begin{quote}
EMSI --- École Marocaine des Sciences de l'Ingénieur Module : Deep
Learning\\
Année universitaire 2025--2026

RAPPORT DE PROJET\\
\textbf{Guide de Référence pour l'Entraînement de Modèles de Deep
Learning}\\
MLP • CNN • RNN • LSTM • GRU • Architectures Hybrides
\end{quote}

\begin{longtable}[]{@{}
  >{\raggedright\arraybackslash}p{(\columnwidth - 2\tabcolsep) * \real{0.5000}}
  >{\raggedright\arraybackslash}p{(\columnwidth - 2\tabcolsep) * \real{0.5000}}@{}}
\toprule()
\begin{minipage}[b]{\linewidth}\raggedright
\begin{quote}
Réalisé par :
\end{quote}
\end{minipage} & \begin{minipage}[b]{\linewidth}\raggedright
\begin{quote}
\textbf{{[}NOM Prénom{]}}
\end{quote}
\end{minipage} \\
\midrule()
\endhead
\begin{minipage}[t]{\linewidth}\raggedright
\begin{quote}
Encadré par :
\end{quote}
\end{minipage} & \begin{minipage}[t]{\linewidth}\raggedright
\begin{quote}
\textbf{Mme. HIDILA Zineb}
\end{quote}
\end{minipage} \\
\begin{minipage}[t]{\linewidth}\raggedright
\begin{quote}
Filière :
\end{quote}
\end{minipage} & \begin{minipage}[t]{\linewidth}\raggedright
\begin{quote}
\textbf{{[}Filière / Année{]}}
\end{quote}
\end{minipage} \\
\bottomrule()
\end{longtable}

\begin{quote}
\emph{Guide de Référence --- Entraînement de Modèles Deep Learning}

\textbf{Table des matières}

Table des matières
...............................................................................................................
2 1. Introduction
.......................................................................................................................
4 1.1 Vue d'ensemble des architectures
...............................................................................
4 2. Préparation des données
..................................................................................................
5 2.1 Exploration et compréhension du dataset
.................................................................... 5
2.2 Prétraitement\\
...............................................................................................................
5 2.2.1 Données tabulaires (MLP)
.....................................................................................
5 2.2.2 Images (CNN)
.......................................................................................................
5 2.2.3 Séquences/Signaux (RNN/LSTM/GRU)
................................................................ 6 2.3
Séparation des données
..............................................................................................
6 3. Architectures en détail\\
.......................................................................................................
7 3.1 Perceptron Multicouche (MLP)
....................................................................................
7 3.1.1 Structure type
........................................................................................................
7 3.1.2 Points
d'attention...................................................................................................
7 3.2 Réseaux de Neurones Convolutifs (CNN)
.................................................................... 8
3.2.1 Structure type (CNN personnalisé)\\
........................................................................
8 3.2.2 Transfert d'apprentissage
......................................................................................
8 3.2.3 Augmentation de données
....................................................................................
8 3.3 Réseaux Récurrents (RNN, LSTM, GRU)
.................................................................... 9
3.3.1 Comparaison des variantes
...................................................................................
9 3.3.2 Configuration type (LSTM)
....................................................................................
9 3.4 Architectures hybrides
...............................................................................................
10 3.4.1 CNN +
LSTM.......................................................................................................
10 3.4.2 CNN + MLP (Fusion multimodale)
.......................................................................
10 3.4.3 Stratégies de fusion
............................................................................................
10 4. Processus d'entraînement
...............................................................................................
11 4.1 Fonctions de perte
.....................................................................................................
11 4.2 Gestion du déséquilibre des classes
..........................................................................
11 4.3 Optimiseurs
...............................................................................................................
11 4.4 Schedulers de learning rate
.......................................................................................
11 4.5 Régularisation
...........................................................................................................
12 5. Guide des hyperparammètres
.........................................................................................
13 5.1 Hyperparammètres clés par architecture
................................................................... 13
5.1.1 MLP
....................................................................................................................
13 5.1.2 CNN
....................................................................................................................
13 5.1.3 RNN / LSTM / GRU
.............................................................................................
13 5.2 Stratégies de recherche d'hyperparammètres
........................................................... 13
\end{quote}

Page 2

\begin{quote}
\emph{Guide de Référence --- Entraînement de Modèles Deep Learning}
\end{quote}

6. Métriques d'évaluation
....................................................................................................
15 6.1 Métriques de classification
.........................................................................................
15 6.2 Visualisations obligatoires
.........................................................................................
15 6.3 Diagnostic par les courbes d'apprentissage
............................................................... 15 7.
Interprétabilité des modèles
............................................................................................
17 8. Bonnes pratiques et erreurs fréquentes
..........................................................................
18 8.1 Checklist avant entraînement
....................................................................................
18 8.2 Erreurs fréquentes
.....................................................................................................
18 8.3 Structure type du code d'entraînement
......................................................................
18 9. Guide de rédaction du rapport
.........................................................................................
20 9.1 Structure attendue
.....................................................................................................
20 9.2 Critères de qualité rédactionnelle
..............................................................................
20 9.3 Ce qu'il faut éviter
......................................................................................................
20 10. Références bibliographiques recommandées
................................................................ 22 10.1
Articles fondateurs
...................................................................................................
22 10.2 Ressources pédagogiques
......................................................................................
22

Page 3

\begin{quote}
\emph{Guide de Référence --- Entraînement de Modèles Deep Learning}

\textbf{1. Introduction}

Ce guide de référence présente la méthodologie complète pour
l'entraînement de modèles de deep learning, couvrant les architectures
fondamentales (MLP, CNN, RNN, LSTM, GRU) ainsi que les architectures
hybrides. Il constitue un cadre méthodologique que l'étudiant doit
adapter à son propre projet.

L'objectif est de fournir une démarche structurée et reproductible,
depuis la préparation des données jusqu'à l'analyse critique des
résultats, en passant par le choix d'architecture, l'optimisation des
hyperparammètres et l'interprétation des modèles.

\textbf{1.1 Vue d'ensemble des architectures}
\end{quote}

\begin{longtable}[]{@{}
  >{\raggedright\arraybackslash}p{(\columnwidth - 6\tabcolsep) * \real{0.2500}}
  >{\raggedright\arraybackslash}p{(\columnwidth - 6\tabcolsep) * \real{0.2500}}
  >{\raggedright\arraybackslash}p{(\columnwidth - 6\tabcolsep) * \real{0.2500}}
  >{\raggedright\arraybackslash}p{(\columnwidth - 6\tabcolsep) * \real{0.2500}}@{}}
\toprule()
\begin{minipage}[b]{\linewidth}\raggedright
\textbf{Architecture}
\end{minipage} & \begin{minipage}[b]{\linewidth}\raggedright
\textbf{Données cibles}
\end{minipage} & \begin{minipage}[b]{\linewidth}\raggedright
\textbf{Biais inductif}
\end{minipage} & \begin{minipage}[b]{\linewidth}\raggedright
\textbf{Cas d'usage typique}
\end{minipage} \\
\midrule()
\endhead
\begin{minipage}[t]{\linewidth}\raggedright
\begin{quote}
MLP

CNN
\end{quote}
\end{minipage} & Tabulaires

Images, grilles 2D & Aucun

Spatial (localité) & \begin{minipage}[t]{\linewidth}\raggedright
\begin{quote}
Classification/régression sur features structurées

Classification d'images, détection d'objets
\end{quote}
\end{minipage} \\
\begin{minipage}[t]{\linewidth}\raggedright
\begin{quote}
RNN

LSTM
\end{quote}
\end{minipage} & \begin{minipage}[t]{\linewidth}\raggedright
\begin{quote}
Séquences courtes

Séquences longues
\end{quote}
\end{minipage} & Temporel (ordre)

Temporel

(mémoire longue) & NLP simple, séries courtes

Traduction, analyse de signaux \\
\bottomrule()
\end{longtable}

\begin{longtable}[]{@{}
  >{\raggedright\arraybackslash}p{(\columnwidth - 6\tabcolsep) * \real{0.2500}}
  >{\raggedright\arraybackslash}p{(\columnwidth - 6\tabcolsep) * \real{0.2500}}
  >{\raggedright\arraybackslash}p{(\columnwidth - 6\tabcolsep) * \real{0.2500}}
  >{\raggedright\arraybackslash}p{(\columnwidth - 6\tabcolsep) * \real{0.2500}}@{}}
\toprule()
\multirow{2}{*}{\begin{minipage}[b]{\linewidth}\raggedright
\begin{quote}
GRU
\end{quote}
\end{minipage}} & \begin{minipage}[b]{\linewidth}\raggedright
Séquences
\end{minipage} & \begin{minipage}[b]{\linewidth}\raggedright
Temporel
\end{minipage} &
\multirow{2}{*}{\begin{minipage}[b]{\linewidth}\raggedright
\begin{quote}
Alternative légère au LSTM
\end{quote}
\end{minipage}} \\
& \begin{minipage}[b]{\linewidth}\raggedright
longues
\end{minipage} & \begin{minipage}[b]{\linewidth}\raggedright
(simplifié)
\end{minipage} \\
\midrule()
\endhead
\begin{minipage}[t]{\linewidth}\raggedright
\begin{quote}
Hybride
\end{quote}
\end{minipage} & Spatio-temporel &
\begin{minipage}[t]{\linewidth}\raggedright
\begin{quote}
Spatial + Temporel
\end{quote}
\end{minipage} & \begin{minipage}[t]{\linewidth}\raggedright
\begin{quote}
Vidéo, signaux multicanaux
\end{quote}
\end{minipage} \\
\bottomrule()
\end{longtable}

\begin{longtable}[]{@{}
  >{\raggedright\arraybackslash}p{(\columnwidth - 0\tabcolsep) * \real{1.0000}}@{}}
\toprule()
\begin{minipage}[b]{\linewidth}\raggedright
\begin{quote}
(CNN+LSTM)
\end{quote}
\end{minipage} \\
\midrule()
\endhead
\bottomrule()
\end{longtable}

Page 4

\begin{quote}
\emph{Guide de Référence --- Entraînement de Modèles Deep Learning}

\textbf{2. Préparation des données}

\textbf{2.1 Exploration et compréhension du dataset}

Avant toute modélisation, une exploration approfondie du dataset est
indispensable. Cette phase doit couvrir les points suivants :

• Description générale : nombre d'échantillons, nombre de
variables/features, type de chaque variable (numérique, catégorielle,
ordinale).

• Distribution de la variable cible : vérifier l'équilibre des classes
et quantifier le ratio de déséquilibre.

• Statistiques descriptives : moyenne, médiane, écart-type, min/max pour
les variables numériques.

• Valeurs manquantes : cartographier les taux de données manquantes par
variable.

• Corrélations : matrice de corrélation pour identifier les redondances
et les relations linéaires.

• Visualisations : histogrammes, boxplots, pairplots, heatmaps de
corrélation.

\textbf{Conseil :} Toujours visualiser les données avant de coder le
modèle. Un simple histogramme peut révéler un déséquilibre critique ou
une distribution inattendue.

\textbf{2.2 Prétraitement}

\textbf{2.2.1 Données tabulaires (MLP)}

1. Gestion des valeurs manquantes : imputation par la médiane
(numérique) ou le mode (catégorielle). Pour les taux supérieurs à 40\%,
envisager la suppression de la variable.

2. Encodage des variables catégorielles : One-Hot Encoding pour les
variables nominales (≤ 10 catégories), Label Encoding pour les
ordinales, Target Encoding si cardinalité élevée.

3. Normalisation/Standardisation : StandardScaler (z-score) ou
MinMaxScaler selon la distribution. Appliquer après le split pour éviter
le data leakage.

4. Sélection de features : suppression des variables à variance nulle,
analyse de corrélation, feature importance préliminaire.

\textbf{Attention :} Toujours fitter le scaler sur le train set
uniquement, puis transformer train, validation et test avec le même
scaler. Fitter sur l'ensemble complet constitue un data leakage.

\textbf{2.2.2 Images (CNN)}

1. Redimensionnement : taille uniforme (ex. 224×224 ou 256×256). Choisir
en fonction de la mémoire GPU disponible.

2. Normalisation : mise à l'échelle {[}0, 1{]}, puis normalisation par
la moyenne et l'écart-type (ImageNet si transfert d'apprentissage :
mean={[}0.485, 0.456, 0.406{]}, std={[}0.229, 0.224, 0.225{]}).

3. Augmentation de données (train uniquement) : rotations,
retournements, recadrage aléatoire, ajustement de luminosité/contraste,
ajout de bruit gaussien.
\end{quote}

Page 5

\begin{quote}
\emph{Guide de Référence --- Entraînement de Modèles Deep Learning}

4. Conversion de format : DICOM vers PNG/JPG si nécessaire, gestion des
canaux (niveaux de gris → 3 canaux pour les modèles pré-entraînés).

\textbf{2.2.3 Séquences/Signaux (RNN/LSTM/GRU)}

1. Rééchantillonnage : uniformiser la fréquence d'échantillonnage (ex.
500 Hz → 100 Hz pour réduire la complexité).
\end{quote}

2. Segmentation : découpage en fenêtres de taille fixe (avec ou sans
recouvrement).

\begin{quote}
3. Normalisation : z-score par canal/variable ou normalisation min-max
par séquence.

4. Padding/Troncature : uniformiser la longueur des séquences (padding
avec zéros ou troncature).

5. Gestion du bruit : filtrage passe-bande si nécessaire, suppression
des artefacts.

\textbf{2.3 Séparation des données}
\end{quote}

\begin{longtable}[]{@{}
  >{\raggedright\arraybackslash}p{(\columnwidth - 6\tabcolsep) * \real{0.2500}}
  >{\raggedright\arraybackslash}p{(\columnwidth - 6\tabcolsep) * \real{0.2500}}
  >{\raggedright\arraybackslash}p{(\columnwidth - 6\tabcolsep) * \real{0.2500}}
  >{\raggedright\arraybackslash}p{(\columnwidth - 6\tabcolsep) * \real{0.2500}}@{}}
\toprule()
\begin{minipage}[b]{\linewidth}\raggedright
\textbf{Ensemble}
\end{minipage} & \begin{minipage}[b]{\linewidth}\raggedright
\textbf{Proportion}
\end{minipage} & \begin{minipage}[b]{\linewidth}\raggedright
\textbf{Rôle}
\end{minipage} & \begin{minipage}[b]{\linewidth}\raggedright
\textbf{Précautions}
\end{minipage} \\
\midrule()
\endhead
\begin{minipage}[t]{\linewidth}\raggedright
\begin{quote}
Train\\
Validation
\end{quote}\strut
\end{minipage} & \begin{minipage}[t]{\linewidth}\raggedright
\begin{quote}
70--80\%\\
10--15\%
\end{quote}\strut
\end{minipage} & \begin{minipage}[t]{\linewidth}\raggedright
Entraînement du modèle\\
Tuning des hyperparammètres\strut
\end{minipage} & \begin{minipage}[t]{\linewidth}\raggedright
\begin{quote}
Stratified split obligatoire\\
Jamais utilisé pour le gradient
\end{quote}\strut
\end{minipage} \\
\begin{minipage}[t]{\linewidth}\raggedright
\begin{quote}
Test
\end{quote}
\end{minipage} & 10--15\% & Évaluation finale unique & Isolé jusqu'à
l'évaluation finale \\
\bottomrule()
\end{longtable}

\begin{quote}
\textbf{Attention :} Le test set ne doit être utilisé qu'une seule fois,
pour l'évaluation finale. Toute utilisation répétée pour ajuster le
modèle constitue un overfitting indirect.
\end{quote}

Page 6

\begin{quote}
\emph{Guide de Référence --- Entraînement de Modèles Deep Learning}

\textbf{3. Architectures en détail}

\textbf{3.1 Perceptron Multicouche (MLP)}

Le MLP est l'architecture la plus directe pour les données tabulaires.
Chaque neurone d'une couche est connecté à tous les neurones de la
couche suivante (fully connected). La sortie d'un neurone est : y =
f(Wᵀx + b).

\textbf{3.1.1 Structure type}
\end{quote}

\begin{longtable}[]{@{}
  >{\raggedright\arraybackslash}p{(\columnwidth - 6\tabcolsep) * \real{0.2500}}
  >{\raggedright\arraybackslash}p{(\columnwidth - 6\tabcolsep) * \real{0.2500}}
  >{\raggedright\arraybackslash}p{(\columnwidth - 6\tabcolsep) * \real{0.2500}}
  >{\raggedright\arraybackslash}p{(\columnwidth - 6\tabcolsep) * \real{0.2500}}@{}}
\toprule()
\begin{minipage}[b]{\linewidth}\raggedright
\textbf{Couche}
\end{minipage} & \begin{minipage}[b]{\linewidth}\raggedright
\textbf{Type}
\end{minipage} & \begin{minipage}[b]{\linewidth}\raggedright
\textbf{Paramètres}
\end{minipage} & \begin{minipage}[b]{\linewidth}\raggedright
\textbf{Remarques}
\end{minipage} \\
\midrule()
\endhead
\begin{minipage}[t]{\linewidth}\raggedright
\begin{quote}
Entrée

Cachée 1
\end{quote}
\end{minipage} & Linear

Linear + ReLU + BN +

Dropout & n\_features → 128

128 → 64 & \begin{minipage}[t]{\linewidth}\raggedright
\begin{quote}
Dimension = nombre de variables

BatchNorm avant ou après ReLU
\end{quote}
\end{minipage} \\
\bottomrule()
\end{longtable}

\begin{longtable}[]{@{}
  >{\raggedright\arraybackslash}p{(\columnwidth - 6\tabcolsep) * \real{0.2500}}
  >{\raggedright\arraybackslash}p{(\columnwidth - 6\tabcolsep) * \real{0.2500}}
  >{\raggedright\arraybackslash}p{(\columnwidth - 6\tabcolsep) * \real{0.2500}}
  >{\raggedright\arraybackslash}p{(\columnwidth - 6\tabcolsep) * \real{0.2500}}@{}}
\toprule()
\begin{minipage}[b]{\linewidth}\raggedright
\begin{quote}
Cachée 2
\end{quote}
\end{minipage} & \begin{minipage}[b]{\linewidth}\raggedright
Linear + ReLU + BN +
\end{minipage} & \begin{minipage}[b]{\linewidth}\raggedright
64 → 32
\end{minipage} & \begin{minipage}[b]{\linewidth}\raggedright
Dropout = 0.2 à 0.5
\end{minipage} \\
\midrule()
\endhead
\bottomrule()
\end{longtable}

Dropout

\begin{longtable}[]{@{}
  >{\raggedright\arraybackslash}p{(\columnwidth - 6\tabcolsep) * \real{0.2500}}
  >{\raggedright\arraybackslash}p{(\columnwidth - 6\tabcolsep) * \real{0.2500}}
  >{\raggedright\arraybackslash}p{(\columnwidth - 6\tabcolsep) * \real{0.2500}}
  >{\raggedright\arraybackslash}p{(\columnwidth - 6\tabcolsep) * \real{0.2500}}@{}}
\toprule()
\begin{minipage}[b]{\linewidth}\raggedright
\begin{quote}
Sortie
\end{quote}
\end{minipage} & \begin{minipage}[b]{\linewidth}\raggedright
\begin{quote}
Linear + Sigmoid/Softmax
\end{quote}
\end{minipage} & \begin{minipage}[b]{\linewidth}\raggedright
32 → n\_classes
\end{minipage} & \begin{minipage}[b]{\linewidth}\raggedright
\begin{quote}
Sigmoid (binaire) / Softmax
\end{quote}
\end{minipage} \\
\midrule()
\endhead
\bottomrule()
\end{longtable}

\begin{longtable}[]{@{}
  >{\raggedright\arraybackslash}p{(\columnwidth - 0\tabcolsep) * \real{1.0000}}@{}}
\toprule()
\begin{minipage}[b]{\linewidth}\raggedright
(multi-classe)
\end{minipage} \\
\midrule()
\endhead
\bottomrule()
\end{longtable}

\begin{quote}
\textbf{3.1.2 Points d'attention}

• Architecture en entonnoir : réduire progressivement la dimension (ex.
256 → 128 → 64 → 32).

• Ne pas aller trop profond : 2 à 4 couches cachées suffisent
généralement pour les données tabulaires.

• Batch Normalization : stabilise et accélère l'entraînement.

• Dropout : régularisation essentielle pour éviter le surapprentissage.
Commencer à 0.3.
\end{quote}

\begin{longtable}[]{@{}
  >{\raggedright\arraybackslash}p{(\columnwidth - 2\tabcolsep) * \real{0.5000}}
  >{\raggedright\arraybackslash}p{(\columnwidth - 2\tabcolsep) * \real{0.5000}}@{}}
\toprule()
\begin{minipage}[b]{\linewidth}\raggedright
•
\end{minipage} & \begin{minipage}[b]{\linewidth}\raggedright
\begin{quote}
Fonction de perte : BCEWithLogitsLoss (binaire) ou CrossEntropyLoss
(multi-classe).
\end{quote}
\end{minipage} \\
\midrule()
\endhead
\bottomrule()
\end{longtable}

Page 7

\begin{quote}
\emph{Guide de Référence --- Entraînement de Modèles Deep Learning}

\textbf{3.2 Réseaux de Neurones Convolutifs (CNN)}

Les CNN exploitent la structure spatiale des images grâce à la
convolution (partage de poids, invariance par translation). Un bloc
convolutif typique comprend : Conv2d → BatchNorm → ReLU → MaxPool.

\textbf{3.2.1 Structure type (CNN personnalisé)}
\end{quote}

\begin{longtable}[]{@{}
  >{\raggedright\arraybackslash}p{(\columnwidth - 6\tabcolsep) * \real{0.2500}}
  >{\raggedright\arraybackslash}p{(\columnwidth - 6\tabcolsep) * \real{0.2500}}
  >{\raggedright\arraybackslash}p{(\columnwidth - 6\tabcolsep) * \real{0.2500}}
  >{\raggedright\arraybackslash}p{(\columnwidth - 6\tabcolsep) * \real{0.2500}}@{}}
\toprule()
\begin{minipage}[b]{\linewidth}\raggedright
\textbf{Bloc}
\end{minipage} & \begin{minipage}[b]{\linewidth}\raggedright
\textbf{Opérations}
\end{minipage} & \begin{minipage}[b]{\linewidth}\raggedright
\textbf{Entrée → Sortie}
\end{minipage} & \begin{minipage}[b]{\linewidth}\raggedright
\textbf{Remarques}
\end{minipage} \\
\midrule()
\endhead
\begin{minipage}[t]{\linewidth}\raggedright
\begin{quote}
Bloc 1

Bloc 2
\end{quote}
\end{minipage} & Conv(3×3, 32) + BN + ReLU + MaxPool(2)

Conv(3×3, 64) + BN + ReLU + MaxPool(2) & 1×256×256 → 32×128×128

32×128×128 → 64×64×64 & Filtres larges en début

Doubler les filtres \\
\begin{minipage}[t]{\linewidth}\raggedright
\begin{quote}
Bloc 3

Bloc 4
\end{quote}
\end{minipage} & \begin{minipage}[t]{\linewidth}\raggedright
\begin{quote}
Conv(3×3, 128) + BN + ReLU + MaxPool(2)

Conv(3×3, 256) + BN + ReLU + AdaptiveAvgPool
\end{quote}
\end{minipage} & 64×64×64 → 128×32×32

128×32×32 → 256×1×1 & Augmenter la profondeur

Global pooling en sortie \\
\begin{minipage}[t]{\linewidth}\raggedright
\begin{quote}
Classif.
\end{quote}
\end{minipage} & \begin{minipage}[t]{\linewidth}\raggedright
\begin{quote}
Linear(256, 128) + ReLU + Dropout + Linear(128, 1)
\end{quote}
\end{minipage} & 256 → 1 & Couches denses finales \\
\bottomrule()
\end{longtable}

\begin{quote}
\textbf{3.2.2 Transfert d'apprentissage}

Le transfert d'apprentissage consiste à réutiliser un modèle
pré-entraîné (ResNet, VGG, EfficientNet) et à adapter les dernières
couches à la tâche cible. Deux stratégies principales :

• Feature extraction : geler toutes les couches convolutives, ne
réentraîner que le classifieur. Rapide, efficace avec peu de données.

• Fine-tuning : dégeler progressivement les dernières couches
convolutives. Plus performant mais nécessite plus de données et un
learning rate faible (1e-4 à 1e-5).

\textbf{Conseil :} Pour le fine-tuning, utiliser un learning rate
différentiel : lr faible pour les couches pré-entraînées (1e-5), lr plus
élevé pour le classifieur (1e-3).

\textbf{3.2.3 Augmentation de données}
\end{quote}

\begin{longtable}[]{@{}
  >{\raggedright\arraybackslash}p{(\columnwidth - 4\tabcolsep) * \real{0.3333}}
  >{\raggedright\arraybackslash}p{(\columnwidth - 4\tabcolsep) * \real{0.3333}}
  >{\raggedright\arraybackslash}p{(\columnwidth - 4\tabcolsep) * \real{0.3333}}@{}}
\toprule()
\begin{minipage}[b]{\linewidth}\raggedright
\textbf{Transformation}
\end{minipage} & \begin{minipage}[b]{\linewidth}\raggedright
\textbf{Paramètres typiques}
\end{minipage} & \begin{minipage}[b]{\linewidth}\raggedright
\textbf{Quand l'utiliser}
\end{minipage} \\
\midrule()
\endhead
\begin{minipage}[t]{\linewidth}\raggedright
\begin{quote}
RandomHorizontalFlip

RandomRotation
\end{quote}
\end{minipage} & p=0.5

degrees=15--30 & \begin{minipage}[t]{\linewidth}\raggedright
\begin{quote}
Toujours (sauf si orientation est discriminante)
\end{quote}

Images sans orientation fixe
\end{minipage} \\
\begin{minipage}[t]{\linewidth}\raggedright
GaussianBlur\\
RandomAffine\strut
\end{minipage} & \begin{minipage}[t]{\linewidth}\raggedright
\begin{quote}
kernel\_size=3\\
translate=(0.1, 0.1)
\end{quote}\strut
\end{minipage} & \begin{minipage}[t]{\linewidth}\raggedright
\begin{quote}
Pour robustesse au bruit d'acquisition Pour simuler des décalages
spatiaux
\end{quote}
\end{minipage} \\
\bottomrule()
\end{longtable}

Page 8

\begin{quote}
\emph{Guide de Référence --- Entraînement de Modèles Deep Learning}

\textbf{3.3 Réseaux Récurrents (RNN, LSTM, GRU)}

Les architectures récurrentes traitent des données séquentielles en
maintenant un état caché qui se propage dans le temps. Le RNN simple
souffre de la disparition du gradient sur les longues séquences, ce qui
a motivé le développement du LSTM et du GRU.

\textbf{3.3.1 Comparaison des variantes}
\end{quote}

\begin{longtable}[]{@{}
  >{\raggedright\arraybackslash}p{(\columnwidth - 6\tabcolsep) * \real{0.2500}}
  >{\raggedright\arraybackslash}p{(\columnwidth - 6\tabcolsep) * \real{0.2500}}
  >{\raggedright\arraybackslash}p{(\columnwidth - 6\tabcolsep) * \real{0.2500}}
  >{\raggedright\arraybackslash}p{(\columnwidth - 6\tabcolsep) * \real{0.2500}}@{}}
\toprule()
\begin{minipage}[b]{\linewidth}\raggedright
\textbf{Critère}
\end{minipage} & \begin{minipage}[b]{\linewidth}\raggedright
\textbf{RNN simple}
\end{minipage} & \begin{minipage}[b]{\linewidth}\raggedright
\textbf{LSTM}
\end{minipage} & \begin{minipage}[b]{\linewidth}\raggedright
\textbf{GRU}
\end{minipage} \\
\midrule()
\endhead
\begin{minipage}[t]{\linewidth}\raggedright
\begin{quote}
État interne\\
Portes
\end{quote}\strut
\end{minipage} & \begin{minipage}[t]{\linewidth}\raggedright
\begin{quote}
1 état caché (h) Aucune
\end{quote}
\end{minipage} & \begin{minipage}[t]{\linewidth}\raggedright
\begin{quote}
2 états (h + cellule c) 3 (entrée, oubli, sortie)
\end{quote}
\end{minipage} & \begin{minipage}[t]{\linewidth}\raggedright
\begin{quote}
1 état caché (h) 2 (reset, update)
\end{quote}
\end{minipage} \\
\begin{minipage}[t]{\linewidth}\raggedright
\begin{quote}
Paramètres\\
Mémoire longue
\end{quote}\strut
\end{minipage} & \begin{minipage}[t]{\linewidth}\raggedright
\begin{quote}
Le moins\\
Faible
\end{quote}\strut
\end{minipage} & \begin{minipage}[t]{\linewidth}\raggedright
\begin{quote}
Le plus\\
Excellente
\end{quote}\strut
\end{minipage} & \begin{minipage}[t]{\linewidth}\raggedright
\begin{quote}
Intermédiaire\\
Très bonne
\end{quote}\strut
\end{minipage} \\
\begin{minipage}[t]{\linewidth}\raggedright
\begin{quote}
Vitesse\\
d'entraînement

Quand l'utiliser
\end{quote}\strut
\end{minipage} & Le plus rapide

Prototypage rapide & Le plus lent

Séquences longues,

défaut & Intermédiaire

Compromis

performance/vitesse \\
\bottomrule()
\end{longtable}

\begin{quote}
\textbf{3.3.2 Configuration type (LSTM)}
\end{quote}

\begin{longtable}[]{@{}
  >{\raggedright\arraybackslash}p{(\columnwidth - 4\tabcolsep) * \real{0.3333}}
  >{\raggedright\arraybackslash}p{(\columnwidth - 4\tabcolsep) * \real{0.3333}}
  >{\raggedright\arraybackslash}p{(\columnwidth - 4\tabcolsep) * \real{0.3333}}@{}}
\toprule()
\begin{minipage}[b]{\linewidth}\raggedright
\textbf{Paramètre}
\end{minipage} & \begin{minipage}[b]{\linewidth}\raggedright
\textbf{Valeur recommandée}
\end{minipage} & \begin{minipage}[b]{\linewidth}\raggedright
\textbf{Remarques}
\end{minipage} \\
\midrule()
\endhead
\begin{minipage}[t]{\linewidth}\raggedright
\begin{quote}
input\_size\\
hidden\_size
\end{quote}\strut
\end{minipage} & Nombre de features par timestep 64 à 256 &
\begin{minipage}[t]{\linewidth}\raggedright
\begin{quote}
Ex. 12 pour ECG 12 dérivations Commencer à 128
\end{quote}
\end{minipage} \\
num\_layers bidirectional & \begin{minipage}[t]{\linewidth}\raggedright
\begin{quote}
1 à 3\\
True
\end{quote}\strut
\end{minipage} & \begin{minipage}[t]{\linewidth}\raggedright
\begin{quote}
2 couches = bon compromis\\
Capture le contexte passé et futur
\end{quote}\strut
\end{minipage} \\
\bottomrule()
\end{longtable}

\begin{longtable}[]{@{}
  >{\raggedright\arraybackslash}p{(\columnwidth - 4\tabcolsep) * \real{0.3333}}
  >{\raggedright\arraybackslash}p{(\columnwidth - 4\tabcolsep) * \real{0.3333}}
  >{\raggedright\arraybackslash}p{(\columnwidth - 4\tabcolsep) * \real{0.3333}}@{}}
\toprule()
\begin{minipage}[b]{\linewidth}\raggedright
\begin{quote}
dropout
\end{quote}
\end{minipage} & \begin{minipage}[b]{\linewidth}\raggedright
0.2 à 0.5
\end{minipage} & \begin{minipage}[b]{\linewidth}\raggedright
\begin{quote}
Entre les couches LSTM (num\_layers \textgreater{} 1)
\end{quote}
\end{minipage} \\
\midrule()
\endhead
\begin{minipage}[t]{\linewidth}\raggedright
\begin{quote}
batch\_first
\end{quote}
\end{minipage} & True & \begin{minipage}[t]{\linewidth}\raggedright
\begin{quote}
Convention PyTorch : (batch, seq, features)
\end{quote}
\end{minipage} \\
\bottomrule()
\end{longtable}

\begin{quote}
\textbf{Note :} Pour la classification, on utilise généralement le
dernier état caché (ou la concaténation forward + backward si
bidirectionnel) passé à une couche dense.
\end{quote}

Page 9

\begin{quote}
\emph{Guide de Référence --- Entraînement de Modèles Deep Learning}

\textbf{3.4 Architectures hybrides}

Les architectures hybrides combinent les forces de plusieurs familles de
réseaux pour traiter des données présentant des structures multiples
(spatiale + temporelle, par exemple).

\textbf{3.4.1 CNN + LSTM}

Cette architecture utilise un CNN comme extracteur de features
spatiales, puis un LSTM pour capturer les dépendances temporelles. Cas
d'usage typiques : classification de vidéos, analyse de signaux
multicanaux avec structure spatiale.

\textbf{Pipeline :} Entrée (batch, seq\_len, C, H, W) → CNN par timestep
→ (batch, seq\_len, n\_features) → LSTM → (batch, hidden) → Dense →
Prédiction.

\textbf{3.4.2 CNN + MLP (Fusion multimodale)}

Pour combiner des données hétérogènes (ex. image + données tabulaires),
on extrait des embeddings de chaque branche puis on les concatène avant
une couche de décision.

\textbf{Pipeline :} Branche image (CNN → embedding) + Branche tabulaire
(MLP → embedding) → Concaténation → Dense → Prédiction.

\textbf{3.4.3 Stratégies de fusion}
\end{quote}

\begin{longtable}[]{@{}
  >{\raggedright\arraybackslash}p{(\columnwidth - 6\tabcolsep) * \real{0.2500}}
  >{\raggedright\arraybackslash}p{(\columnwidth - 6\tabcolsep) * \real{0.2500}}
  >{\raggedright\arraybackslash}p{(\columnwidth - 6\tabcolsep) * \real{0.2500}}
  >{\raggedright\arraybackslash}p{(\columnwidth - 6\tabcolsep) * \real{0.2500}}@{}}
\toprule()
\begin{minipage}[b]{\linewidth}\raggedright
\textbf{Stratégie}
\end{minipage} & \begin{minipage}[b]{\linewidth}\raggedright
\textbf{Description}
\end{minipage} & \begin{minipage}[b]{\linewidth}\raggedright
\textbf{Avantages}
\end{minipage} & \begin{minipage}[b]{\linewidth}\raggedright
\textbf{Inconvénients}
\end{minipage} \\
\midrule()
\endhead
\begin{minipage}[t]{\linewidth}\raggedright
\begin{quote}
Early fusion

Late fusion
\end{quote}
\end{minipage} & \begin{minipage}[t]{\linewidth}\raggedright
Concaténation des données brutes en entrée

\begin{quote}
Concaténation des embeddings avant la décision
\end{quote}
\end{minipage} & Simple à implémenter

Flexible, modulaire & \begin{minipage}[t]{\linewidth}\raggedright
\begin{quote}
Dimensions incompatibles
\end{quote}

Interactions limitées
\end{minipage} \\
\bottomrule()
\end{longtable}

\begin{longtable}[]{@{}
  >{\raggedright\arraybackslash}p{(\columnwidth - 6\tabcolsep) * \real{0.2500}}
  >{\raggedright\arraybackslash}p{(\columnwidth - 6\tabcolsep) * \real{0.2500}}
  >{\raggedright\arraybackslash}p{(\columnwidth - 6\tabcolsep) * \real{0.2500}}
  >{\raggedright\arraybackslash}p{(\columnwidth - 6\tabcolsep) * \real{0.2500}}@{}}
\toprule()
\multirow{2}{*}{\begin{minipage}[b]{\linewidth}\raggedright
\begin{quote}
Intermediate fusion
\end{quote}
\end{minipage}} & \begin{minipage}[b]{\linewidth}\raggedright
\begin{quote}
Fusion à un niveau
\end{quote}
\end{minipage} &
\multirow{2}{*}{\begin{minipage}[b]{\linewidth}\raggedright
Interactions riches
\end{minipage}} & \begin{minipage}[b]{\linewidth}\raggedright
Plus complexe à
\end{minipage} \\
& \begin{minipage}[b]{\linewidth}\raggedright
\begin{quote}
intermédiaire du réseau
\end{quote}
\end{minipage} & & \begin{minipage}[b]{\linewidth}\raggedright
concevoir
\end{minipage} \\
\midrule()
\endhead
\multirow{2}{*}{\begin{minipage}[t]{\linewidth}\raggedright
\begin{quote}
Attention-based
\end{quote}
\end{minipage}} & \begin{minipage}[t]{\linewidth}\raggedright
\begin{quote}
Mécanisme d'attention
\end{quote}
\end{minipage} & \multirow{2}{*}{Pondération adaptative} & Coût
computationnel \\
& cross-modal & & élevé \\
\bottomrule()
\end{longtable}

Page 10

\begin{quote}
\emph{Guide de Référence --- Entraînement de Modèles Deep Learning}

\textbf{4. Processus d'entraînement}

\textbf{4.1 Fonctions de perte}
\end{quote}

\begin{longtable}[]{@{}
  >{\raggedright\arraybackslash}p{(\columnwidth - 6\tabcolsep) * \real{0.2500}}
  >{\raggedright\arraybackslash}p{(\columnwidth - 6\tabcolsep) * \real{0.2500}}
  >{\raggedright\arraybackslash}p{(\columnwidth - 6\tabcolsep) * \real{0.2500}}
  >{\raggedright\arraybackslash}p{(\columnwidth - 6\tabcolsep) * \real{0.2500}}@{}}
\toprule()
\begin{minipage}[b]{\linewidth}\raggedright
\textbf{Tâche}
\end{minipage} & \begin{minipage}[b]{\linewidth}\raggedright
\begin{quote}
\textbf{Fonction de perte PyTorch}
\end{quote}
\end{minipage} & \begin{minipage}[b]{\linewidth}\raggedright
\textbf{Activation de sortie}
\end{minipage} & \begin{minipage}[b]{\linewidth}\raggedright
\textbf{Quand l'utiliser}
\end{minipage} \\
\midrule()
\endhead
\begin{minipage}[t]{\linewidth}\raggedright
\begin{quote}
Classification binaire Classification multi-classe
\end{quote}
\end{minipage} & \begin{minipage}[t]{\linewidth}\raggedright
\begin{quote}
nn.BCEWithLogitsLoss() nn.CrossEntropyLoss()
\end{quote}
\end{minipage} & \begin{minipage}[t]{\linewidth}\raggedright
\begin{quote}
Aucune (logits) Aucune (logits)
\end{quote}
\end{minipage} & \begin{minipage}[t]{\linewidth}\raggedright
\begin{quote}
2 classes, label unique N classes mutuellement exclusives
\end{quote}
\end{minipage} \\
\begin{minipage}[t]{\linewidth}\raggedright
\begin{quote}
Classification multi-labels

Régression
\end{quote}
\end{minipage} & nn.BCEWithLogitsLoss()

nn.MSELoss() ou

nn.L1Loss() & Aucune (logits)

Aucune (linéaire) & \begin{minipage}[t]{\linewidth}\raggedright
\begin{quote}
Plusieurs labels simultanés

Prédiction de valeurs continues
\end{quote}
\end{minipage} \\
\bottomrule()
\end{longtable}

\begin{quote}
\textbf{Attention :} BCEWithLogitsLoss attend des logits (pas de
sigmoïde en sortie).

CrossEntropyLoss attend des logits (pas de softmax). Appliquer
sigmoïde/softmax avant ces loss fonctions est une erreur fréquente.

\textbf{4.2 Gestion du déséquilibre des classes}

Le déséquilibre des classes est un problème récurrent en deep learning
appliqué. Plusieurs stratégies peuvent être combinées :

• Pondération de la loss : calculer les poids inversement proportionnels
à la fréquence de chaque classe (paramètre weight de CrossEntropyLoss ou
pos\_weight de BCEWithLogitsLoss).

• Suréchantillonnage de la classe minoritaire : SMOTE (données
tabulaires), augmentation ciblée (images).

• Sous-échantillonnage de la classe majoritaire : à utiliser avec
précaution (perte d'information).

• Focal Loss : variante de la BCE qui réduit le poids des exemples
faciles. Utile pour les déséquilibres sévères.

\textbf{4.3 Optimiseurs}
\end{quote}

\begin{longtable}[]{@{}
  >{\raggedright\arraybackslash}p{(\columnwidth - 6\tabcolsep) * \real{0.2500}}
  >{\raggedright\arraybackslash}p{(\columnwidth - 6\tabcolsep) * \real{0.2500}}
  >{\raggedright\arraybackslash}p{(\columnwidth - 6\tabcolsep) * \real{0.2500}}
  >{\raggedright\arraybackslash}p{(\columnwidth - 6\tabcolsep) * \real{0.2500}}@{}}
\toprule()
\begin{minipage}[b]{\linewidth}\raggedright
\textbf{Optimiseur}
\end{minipage} & \begin{minipage}[b]{\linewidth}\raggedright
\textbf{Learning rate typique}
\end{minipage} & \begin{minipage}[b]{\linewidth}\raggedright
\textbf{Avantages}
\end{minipage} & \begin{minipage}[b]{\linewidth}\raggedright
\textbf{Cas d'usage}
\end{minipage} \\
\midrule()
\endhead
\begin{minipage}[t]{\linewidth}\raggedright
\begin{quote}
SGD + Momentum

Adam
\end{quote}
\end{minipage} & 0.01 à 0.1

1e-3 à 1e-4 & Bonne généralisation

Convergence rapide,

adaptatif & \begin{minipage}[t]{\linewidth}\raggedright
\begin{quote}
CNN avec transfert, entraînement long

Défaut pour la plupart des cas
\end{quote}
\end{minipage} \\
\begin{minipage}[t]{\linewidth}\raggedright
\begin{quote}
AdamW

RMSProp
\end{quote}
\end{minipage} & 1e-3 à 1e-4

1e-3 & \begin{minipage}[t]{\linewidth}\raggedright
\begin{quote}
Adam + weight decay découplé
\end{quote}

Bon pour les RNN
\end{minipage} & \begin{minipage}[t]{\linewidth}\raggedright
Fine-tuning de modèles pré-entraînés

\begin{quote}
Séquences, gradients instables
\end{quote}
\end{minipage} \\
\bottomrule()
\end{longtable}

\begin{quote}
\textbf{4.4 Schedulers de learning rate}
\end{quote}

Page 11

\begin{quote}
\emph{Guide de Référence --- Entraînement de Modèles Deep Learning}

Le scheduler ajuste le learning rate au cours de l'entraînement pour
améliorer la

convergence :

• StepLR : réduit le lr d'un facteur gamma tous les step\_size époques.
Simple et

efficace.

• ReduceLROnPlateau : réduit le lr quand la métrique de validation
stagne.

Recommandé comme défaut.

• CosineAnnealingLR : décroissance en cosinus. Populaire pour le
fine-tuning.

• OneCycleLR : monte puis descend le lr en un cycle. Peut accélérer la
convergence.

\textbf{4.5 Régularisation}
\end{quote}

\begin{longtable}[]{@{}
  >{\raggedright\arraybackslash}p{(\columnwidth - 6\tabcolsep) * \real{0.2500}}
  >{\raggedright\arraybackslash}p{(\columnwidth - 6\tabcolsep) * \real{0.2500}}
  >{\raggedright\arraybackslash}p{(\columnwidth - 6\tabcolsep) * \real{0.2500}}
  >{\raggedright\arraybackslash}p{(\columnwidth - 6\tabcolsep) * \real{0.2500}}@{}}
\toprule()
\begin{minipage}[b]{\linewidth}\raggedright
\textbf{Technique}
\end{minipage} & \begin{minipage}[b]{\linewidth}\raggedright
\textbf{Où l'appliquer}
\end{minipage} & \begin{minipage}[b]{\linewidth}\raggedright
\textbf{Paramètres}
\end{minipage} & \begin{minipage}[b]{\linewidth}\raggedright
\textbf{Effet}
\end{minipage} \\
\midrule()
\endhead
\begin{minipage}[t]{\linewidth}\raggedright
\begin{quote}
Dropout

Batch Normalization
\end{quote}
\end{minipage} & \begin{minipage}[t]{\linewidth}\raggedright
\begin{quote}
Après les couches denses/LSTM

Après Conv/Linear, avant ou après ReLU
\end{quote}
\end{minipage} & p = 0.2 à 0.5

momentum=0.1 & \begin{minipage}[t]{\linewidth}\raggedright
\begin{quote}
Désactive des neurones aléatoirement

Normalise les activations par batch
\end{quote}
\end{minipage} \\
\begin{minipage}[t]{\linewidth}\raggedright
\begin{quote}
Weight Decay (L2) Early Stopping
\end{quote}
\end{minipage} & \begin{minipage}[t]{\linewidth}\raggedright
\begin{quote}
Via l'optimiseur\\
Boucle d'entraînement
\end{quote}\strut
\end{minipage} & \begin{minipage}[t]{\linewidth}\raggedright
\begin{quote}
1e-4 à 1e-2\\
patience = 5 à 15
\end{quote}\strut
\end{minipage} & \begin{minipage}[t]{\linewidth}\raggedright
Pénalise les grands poids Arrête quand val\_loss\\
stagne\strut
\end{minipage} \\
\begin{minipage}[t]{\linewidth}\raggedright
\begin{quote}
Data Augmentation

Label Smoothing
\end{quote}
\end{minipage} & Pipeline de données (train)

Fonction de perte & Variable

ε = 0.1 & \begin{minipage}[t]{\linewidth}\raggedright
\begin{quote}
Augmente la diversité virtuelle

Adoucit les cibles (moins de confiance)
\end{quote}
\end{minipage} \\
\bottomrule()
\end{longtable}

Page 12

\begin{quote}
\emph{Guide de Référence --- Entraînement de Modèles Deep Learning}

\textbf{5. Guide des hyperparammètres}

\textbf{5.1 Hyperparammètres clés par architecture}

\textbf{5.1.1 MLP}
\end{quote}

\begin{longtable}[]{@{}
  >{\raggedright\arraybackslash}p{(\columnwidth - 4\tabcolsep) * \real{0.3333}}
  >{\raggedright\arraybackslash}p{(\columnwidth - 4\tabcolsep) * \real{0.3333}}
  >{\raggedright\arraybackslash}p{(\columnwidth - 4\tabcolsep) * \real{0.3333}}@{}}
\toprule()
\begin{minipage}[b]{\linewidth}\raggedright
\textbf{Hyperparammètre}
\end{minipage} & \begin{minipage}[b]{\linewidth}\raggedright
\textbf{Plage de recherche}
\end{minipage} & \begin{minipage}[b]{\linewidth}\raggedright
\textbf{Défaut recommandé}
\end{minipage} \\
\midrule()
\endhead
\begin{minipage}[t]{\linewidth}\raggedright
\begin{quote}
Nombre de couches cachées Neurones par couche
\end{quote}
\end{minipage} & \begin{minipage}[t]{\linewidth}\raggedright
\begin{quote}
2 à 4\\
32 à 512
\end{quote}\strut
\end{minipage} & \begin{minipage}[t]{\linewidth}\raggedright
\begin{quote}
3\\
128, 64, 32 (entonnoir)
\end{quote}\strut
\end{minipage} \\
\begin{minipage}[t]{\linewidth}\raggedright
\begin{quote}
Learning rate\\
Batch size
\end{quote}\strut
\end{minipage} & \begin{minipage}[t]{\linewidth}\raggedright
\begin{quote}
1e-4 à 1e-2\\
32 à 256
\end{quote}\strut
\end{minipage} & \begin{minipage}[t]{\linewidth}\raggedright
\begin{quote}
1e-3 (Adam)\\
64
\end{quote}\strut
\end{minipage} \\
\begin{minipage}[t]{\linewidth}\raggedright
Dropout\\
Époques\strut
\end{minipage} & \begin{minipage}[t]{\linewidth}\raggedright
\begin{quote}
0.1 à 0.5\\
50 à 200
\end{quote}\strut
\end{minipage} & \begin{minipage}[t]{\linewidth}\raggedright
\begin{quote}
0.3\\
100 + early stopping
\end{quote}\strut
\end{minipage} \\
\begin{minipage}[t]{\linewidth}\raggedright
\begin{quote}
Weight decay
\end{quote}
\end{minipage} & 0 à 1e-2 & 1e-4 \\
\bottomrule()
\end{longtable}

\begin{quote}
\textbf{5.1.2 CNN}
\end{quote}

\begin{longtable}[]{@{}
  >{\raggedright\arraybackslash}p{(\columnwidth - 4\tabcolsep) * \real{0.3333}}
  >{\raggedright\arraybackslash}p{(\columnwidth - 4\tabcolsep) * \real{0.3333}}
  >{\raggedright\arraybackslash}p{(\columnwidth - 4\tabcolsep) * \real{0.3333}}@{}}
\toprule()
\begin{minipage}[b]{\linewidth}\raggedright
\textbf{Hyperparammètre}
\end{minipage} & \begin{minipage}[b]{\linewidth}\raggedright
\textbf{Plage de recherche}
\end{minipage} & \begin{minipage}[b]{\linewidth}\raggedright
\textbf{Défaut recommandé}
\end{minipage} \\
\midrule()
\endhead
\begin{minipage}[t]{\linewidth}\raggedright
\begin{quote}
Nombre de blocs convolutifs Filtres par bloc
\end{quote}
\end{minipage} & \begin{minipage}[t]{\linewidth}\raggedright
\begin{quote}
3 à 6\\
16 à 512
\end{quote}\strut
\end{minipage} & \begin{minipage}[t]{\linewidth}\raggedright
\begin{quote}
4\\
32, 64, 128, 256 (doublement)
\end{quote}\strut
\end{minipage} \\
\begin{minipage}[t]{\linewidth}\raggedright
\begin{quote}
Taille du noyau\\
Learning rate (from scratch)
\end{quote}\strut
\end{minipage} & \begin{minipage}[t]{\linewidth}\raggedright
\begin{quote}
3×3, 5×5, 7×7\\
1e-3 à 1e-2
\end{quote}\strut
\end{minipage} & \begin{minipage}[t]{\linewidth}\raggedright
\begin{quote}
3×3 (standard)\\
1e-3
\end{quote}\strut
\end{minipage} \\
\begin{minipage}[t]{\linewidth}\raggedright
\begin{quote}
Époques\\
Augmentation
\end{quote}\strut
\end{minipage} & \begin{minipage}[t]{\linewidth}\raggedright
\begin{quote}
20 à 100\\
Légère à agressive
\end{quote}\strut
\end{minipage} & \begin{minipage}[t]{\linewidth}\raggedright
\begin{quote}
50 + early stopping Modérée
\end{quote}
\end{minipage} \\
\bottomrule()
\end{longtable}

\begin{quote}
\textbf{5.1.3 RNN / LSTM / GRU}
\end{quote}

\begin{longtable}[]{@{}
  >{\raggedright\arraybackslash}p{(\columnwidth - 4\tabcolsep) * \real{0.3333}}
  >{\raggedright\arraybackslash}p{(\columnwidth - 4\tabcolsep) * \real{0.3333}}
  >{\raggedright\arraybackslash}p{(\columnwidth - 4\tabcolsep) * \real{0.3333}}@{}}
\toprule()
\begin{minipage}[b]{\linewidth}\raggedright
\textbf{Hyperparammètre}
\end{minipage} & \begin{minipage}[b]{\linewidth}\raggedright
\textbf{Plage de recherche}
\end{minipage} & \begin{minipage}[b]{\linewidth}\raggedright
\textbf{Défaut recommandé}
\end{minipage} \\
\midrule()
\endhead
\begin{minipage}[t]{\linewidth}\raggedright
\begin{quote}
Bidirectionnel\\
Dropout inter-couches
\end{quote}\strut
\end{minipage} & \begin{minipage}[t]{\linewidth}\raggedright
\begin{quote}
True / False\\
0.1 à 0.5
\end{quote}\strut
\end{minipage} & \begin{minipage}[t]{\linewidth}\raggedright
\begin{quote}
True\\
0.3
\end{quote}\strut
\end{minipage} \\
\begin{minipage}[t]{\linewidth}\raggedright
\begin{quote}
Learning rate\\
Gradient clipping
\end{quote}\strut
\end{minipage} & \begin{minipage}[t]{\linewidth}\raggedright
\begin{quote}
1e-4 à 1e-2\\
0.5 à 5.0
\end{quote}\strut
\end{minipage} & \begin{minipage}[t]{\linewidth}\raggedright
\begin{quote}
1e-3\\
1.0
\end{quote}\strut
\end{minipage} \\
\begin{minipage}[t]{\linewidth}\raggedright
\begin{quote}
Longueur de séquence
\end{quote}
\end{minipage} & Dépend du signal & 1000 (après rééchantillonnage) \\
\bottomrule()
\end{longtable}

\begin{quote}
\textbf{5.2 Stratégies de recherche d'hyperparammètres}

• Grid Search : exhaustif mais coûteux. Réserver pour 2--3
hyperparammètres avec peu de valeurs.

• Random Search : plus efficace que le grid search pour des espaces
larges. 50--100 itérations.
\end{quote}

Page 13

\begin{quote}
\emph{Guide de Référence --- Entraînement de Modèles Deep Learning}
\end{quote}

\begin{longtable}[]{@{}
  >{\raggedright\arraybackslash}p{(\columnwidth - 2\tabcolsep) * \real{0.5000}}
  >{\raggedright\arraybackslash}p{(\columnwidth - 2\tabcolsep) * \real{0.5000}}@{}}
\toprule()
\begin{minipage}[b]{\linewidth}\raggedright
•
\end{minipage} & \begin{minipage}[b]{\linewidth}\raggedright
\begin{quote}
Recherche manuelle guidée : commencer par les défauts, varier un
paramètre à la
\end{quote}
\end{minipage} \\
\midrule()
\endhead
\bottomrule()
\end{longtable}

\begin{quote}
fois. Approche pragmatique recommandée pour les projets étudiants.

\textbf{Conseil :} Commencer par le learning rate (le plus impactant),
puis batch size, puis profondeur, puis régularisation. Toujours fixer un
seed pour la reproductibilité.
\end{quote}

Page 14

\begin{quote}
\emph{Guide de Référence --- Entraînement de Modèles Deep Learning}

\textbf{6. Métriques d'évaluation}

\textbf{6.1 Métriques de classification}
\end{quote}

\begin{longtable}[]{@{}
  >{\raggedright\arraybackslash}p{(\columnwidth - 4\tabcolsep) * \real{0.3333}}
  >{\raggedright\arraybackslash}p{(\columnwidth - 4\tabcolsep) * \real{0.3333}}
  >{\raggedright\arraybackslash}p{(\columnwidth - 4\tabcolsep) * \real{0.3333}}@{}}
\toprule()
\begin{minipage}[b]{\linewidth}\raggedright
\textbf{Métrique}
\end{minipage} & \begin{minipage}[b]{\linewidth}\raggedright
\textbf{Formule / Description}
\end{minipage} & \begin{minipage}[b]{\linewidth}\raggedright
\textbf{Quand l'utiliser}
\end{minipage} \\
\midrule()
\endhead
\begin{minipage}[t]{\linewidth}\raggedright
\begin{quote}
Accuracy

Précision
\end{quote}
\end{minipage} & TP + TN / Total

TP / (TP + FP) & \begin{minipage}[t]{\linewidth}\raggedright
\begin{quote}
Classes équilibrées uniquement
\end{quote}

Coût élevé des faux positifs
\end{minipage} \\
\begin{minipage}[t]{\linewidth}\raggedright
\begin{quote}
Rappel (Sensibilité)

F1-Score
\end{quote}
\end{minipage} & TP / (TP + FN)

2 × (Précision × Rappel) / (Précision +

Rappel) & \begin{minipage}[t]{\linewidth}\raggedright
\begin{quote}
Coût élevé des faux négatifs (médical)
\end{quote}

Compromis précision/rappel
\end{minipage} \\
\begin{minipage}[t]{\linewidth}\raggedright
\begin{quote}
AUC-ROC

Matrice de confusion
\end{quote}
\end{minipage} & Aire sous la courbe ROC

Tableau TP/TN/FP/FN & \begin{minipage}[t]{\linewidth}\raggedright
\begin{quote}
Métrique globale de discrimination
\end{quote}

Analyse détaillée des erreurs
\end{minipage} \\
\begin{minipage}[t]{\linewidth}\raggedright
\begin{quote}
Macro F1

Spécificité
\end{quote}
\end{minipage} & Moyenne des F1 par classe

TN / (TN + FP) & \begin{minipage}[t]{\linewidth}\raggedright
\begin{quote}
Classes multiples déséquilibrées
\end{quote}

Taux de vrais négatifs
\end{minipage} \\
\bottomrule()
\end{longtable}

\begin{quote}
\textbf{Attention :} En médecine, le rappel (sensibilité) est souvent
plus critique que la précision. Un faux négatif (maladie non détectée)
est généralement plus dangereux qu'un faux positif (examen
supplémentaire inutile).

\textbf{6.2 Visualisations obligatoires}\\
Chaque expérience doit produire les visualisations suivantes :

1. Courbes d'apprentissage : loss et métrique principale (accuracy, F1)
sur train et validation en fonction des époques. Permet de diagnostiquer
l'overfitting/underfitting.

2. Matrice de confusion : sur le test set, normalisée ou non. Identifie
les confusions inter-classes.

3. Courbe ROC et AUC : pour chaque classe en classification
binaire/multi-classe. 4. Tableau récapitulatif : toutes les
configurations testées avec leurs hyperparammètres et métriques.

\textbf{6.3 Diagnostic par les courbes d'apprentissage}
\end{quote}

\begin{longtable}[]{@{}
  >{\raggedright\arraybackslash}p{(\columnwidth - 4\tabcolsep) * \real{0.3333}}
  >{\raggedright\arraybackslash}p{(\columnwidth - 4\tabcolsep) * \real{0.3333}}
  >{\raggedright\arraybackslash}p{(\columnwidth - 4\tabcolsep) * \real{0.3333}}@{}}
\toprule()
\begin{minipage}[b]{\linewidth}\raggedright
\textbf{Observation}
\end{minipage} & \begin{minipage}[b]{\linewidth}\raggedright
\textbf{Diagnostic}
\end{minipage} & \begin{minipage}[b]{\linewidth}\raggedright
\textbf{Action corrective}
\end{minipage} \\
\midrule()
\endhead
\begin{minipage}[t]{\linewidth}\raggedright
\begin{quote}
train\_loss ↓↓, val\_loss ↓ puis ↑

train\_loss ↓ lent, val\_loss ↓ lent
\end{quote}
\end{minipage} & Overfitting

Underfitting & \begin{minipage}[t]{\linewidth}\raggedright
\begin{quote}
Augmenter dropout, réduire modèle, early stopping

Augmenter le modèle, réduire régularisation, augmenter lr
\end{quote}
\end{minipage} \\
\begin{minipage}[t]{\linewidth}\raggedright
\begin{quote}
train\_loss et val\_loss stagnent haut

train\_loss et val\_loss ↓ ensemble
\end{quote}
\end{minipage} & \begin{minipage}[t]{\linewidth}\raggedright
\begin{quote}
Modèle trop simple ou lr trop faible
\end{quote}

Bon entraînement
\end{minipage} & Augmenter capacité ou lr

Continuer, surveiller la divergence \\
\begin{minipage}[t]{\linewidth}\raggedright
\begin{quote}
val\_loss oscille fortement
\end{quote}
\end{minipage} & \begin{minipage}[t]{\linewidth}\raggedright
\begin{quote}
Learning rate trop élevé ou batch trop petit
\end{quote}
\end{minipage} & Réduire lr, augmenter batch size \\
\bottomrule()
\end{longtable}

Page 15

\begin{quote}
\emph{Guide de Référence --- Entraînement de Modèles Deep Learning}
\end{quote}

Page 16

\begin{quote}
\emph{Guide de Référence --- Entraînement de Modèles Deep Learning}

\textbf{7. Interprétabilité des modèles}

L'interprétabilité est un enjeu majeur en deep learning appliqué,
particulièrement dans les
\end{quote}

domaines critiques (santé, finance, sécurité). Un modèle performant mais
opaque est

\begin{quote}
difficilement adoptable en pratique clinique.
\end{quote}

\begin{longtable}[]{@{}
  >{\raggedright\arraybackslash}p{(\columnwidth - 4\tabcolsep) * \real{0.3333}}
  >{\raggedright\arraybackslash}p{(\columnwidth - 4\tabcolsep) * \real{0.3333}}
  >{\raggedright\arraybackslash}p{(\columnwidth - 4\tabcolsep) * \real{0.3333}}@{}}
\toprule()
\begin{minipage}[b]{\linewidth}\raggedright
\textbf{Architecture}
\end{minipage} & \begin{minipage}[b]{\linewidth}\raggedright
\textbf{Méthode d'interprétation}
\end{minipage} & \begin{minipage}[b]{\linewidth}\raggedright
\textbf{Ce qu'elle révèle}
\end{minipage} \\
\midrule()
\endhead
\begin{minipage}[t]{\linewidth}\raggedright
\begin{quote}
CNN\\
CNN
\end{quote}\strut
\end{minipage} & \begin{minipage}[t]{\linewidth}\raggedright
\begin{quote}
Grad-CAM\\
Saliency Maps
\end{quote}\strut
\end{minipage} & \begin{minipage}[t]{\linewidth}\raggedright
\begin{quote}
Quelles régions spatiales activent la décision Gradients de la sortie
par rapport aux pixels d'entrée
\end{quote}
\end{minipage} \\
\begin{minipage}[t]{\linewidth}\raggedright
\begin{quote}
RNN/LSTM Tous
\end{quote}
\end{minipage} & \begin{minipage}[t]{\linewidth}\raggedright
Gradient par rapport à l'entrée t-SNE / UMAP des\\
embeddings\strut
\end{minipage} & \begin{minipage}[t]{\linewidth}\raggedright
\begin{quote}
Quels segments du signal sont discriminants Visualisation de l'espace
latent appris
\end{quote}
\end{minipage} \\
\bottomrule()
\end{longtable}

\begin{quote}
\textbf{Conseil :} Inclure au moins une méthode d'interprétation par
architecture dans le rapport.

Discuter la cohérence des résultats avec les connaissances du domaine.
\end{quote}

Page 17

\begin{quote}
\emph{Guide de Référence --- Entraînement de Modèles Deep Learning}

\textbf{8. Bonnes pratiques et erreurs fréquentes}

\textbf{8.1 Checklist avant entraînement}

1. Seed fixé (torch.manual\_seed, np.random.seed, random.seed) pour la
reproductibilité.

2. Data leakage vérifié : aucune fuite d'information du test vers le
train.

3. Dimensions vérifiées : input shape, output shape, nombre de classes.

4. Baseline établie : modèle simple (régression logistique, prédiction
majoritaire) comme référence.

5. GPU disponible : torch.cuda.is\_available(), device =
torch.device('cuda' if available).

6. DataLoader configuré : shuffle=True pour train, shuffle=False pour
val/test, num\_workers adapté.

\textbf{8.2 Erreurs fréquentes}
\end{quote}

\begin{longtable}[]{@{}
  >{\raggedright\arraybackslash}p{(\columnwidth - 4\tabcolsep) * \real{0.3333}}
  >{\raggedright\arraybackslash}p{(\columnwidth - 4\tabcolsep) * \real{0.3333}}
  >{\raggedright\arraybackslash}p{(\columnwidth - 4\tabcolsep) * \real{0.3333}}@{}}
\toprule()
\begin{minipage}[b]{\linewidth}\raggedright
\textbf{Erreur}
\end{minipage} & \begin{minipage}[b]{\linewidth}\raggedright
\textbf{Conséquence}
\end{minipage} & \begin{minipage}[b]{\linewidth}\raggedright
\textbf{Solution}
\end{minipage} \\
\midrule()
\endhead
\begin{minipage}[t]{\linewidth}\raggedright
\begin{quote}
Oublier model.eval() en évaluation

Softmax avant\\
CrossEntropyLoss
\end{quote}\strut
\end{minipage} & \begin{minipage}[t]{\linewidth}\raggedright
\begin{quote}
Dropout et BN actifs → métriques faussées

Double application → convergence dégradée
\end{quote}
\end{minipage} & Toujours model.eval() + torch.no\_grad()

CrossEntropyLoss attend des logits bruts \\
\begin{minipage}[t]{\linewidth}\raggedright
\begin{quote}
Normaliser sur tout le dataset

Ignorer le déséquilibre des classes
\end{quote}
\end{minipage} & \begin{minipage}[t]{\linewidth}\raggedright
Data leakage → métriques optimistes

\begin{quote}
Modèle biaisé vers la classe majoritaire
\end{quote}
\end{minipage} & Fitter le scaler sur train uniquement

Pondération de la loss, SMOTE, augmentation \\
\begin{minipage}[t]{\linewidth}\raggedright
\begin{quote}
Batch size trop grand\\
Pas de gradient clipping (RNN)
\end{quote}\strut
\end{minipage} & \begin{minipage}[t]{\linewidth}\raggedright
\begin{quote}
Mauvaise généralisation Explosion du gradient → NaN
\end{quote}
\end{minipage} & \begin{minipage}[t]{\linewidth}\raggedright
32 à 64 par défaut\\
torch.nn.utils.clip\_grad\_norm\_(max\_norm=1.0)\strut
\end{minipage} \\
\begin{minipage}[t]{\linewidth}\raggedright
\begin{quote}
Oublier\\
optimizer.zero\_grad()

Pas d'early stopping
\end{quote}\strut
\end{minipage} & \begin{minipage}[t]{\linewidth}\raggedright
Accumulation de gradients → instabilité

\begin{quote}
Surapprentissage non détecté
\end{quote}
\end{minipage} & Appeler en début de chaque itération

Monitorer val\_loss, patience=5--15 \\
\bottomrule()
\end{longtable}

\begin{quote}
\textbf{8.3 Structure type du code d'entraînement}

Chaque notebook ou script d'entraînement doit suivre cette structure
logique :

1. Imports et configuration (device, seeds, hyperparammètres).

2. Chargement et exploration des données.

3. Prétraitement et création des DataLoaders.

4. Définition du modèle (classe héritant de nn.Module).

5. Définition de la loss, de l'optimiseur et du scheduler.

6. Boucle d'entraînement avec logging (train\_loss, val\_loss, métriques
par époque).

7. Évaluation sur le test set (une seule fois, après sélection du
meilleur modèle). 8. Visualisations (courbes, matrices, Grad-CAM, etc.).

9. Analyse critique et conclusions.
\end{quote}

Page 18

\begin{quote}
\emph{Guide de Référence --- Entraînement de Modèles Deep Learning}
\end{quote}

Page 19

\begin{quote}
\emph{Guide de Référence --- Entraînement de Modèles Deep Learning}

\textbf{9. Guide de rédaction du rapport}

Le rapport constitue le livrable principal du projet. Il doit démontrer
la compréhension

théorique, la rigueur expérimentale et la capacité d'analyse critique de
l'étudiant.

\textbf{9.1 Structure attendue}
\end{quote}

\begin{longtable}[]{@{}
  >{\raggedright\arraybackslash}p{(\columnwidth - 4\tabcolsep) * \real{0.3333}}
  >{\raggedright\arraybackslash}p{(\columnwidth - 4\tabcolsep) * \real{0.3333}}
  >{\raggedright\arraybackslash}p{(\columnwidth - 4\tabcolsep) * \real{0.3333}}@{}}
\toprule()
\begin{minipage}[b]{\linewidth}\raggedright
\textbf{Section}
\end{minipage} & \begin{minipage}[b]{\linewidth}\raggedright
\textbf{Contenu attendu}
\end{minipage} & \begin{minipage}[b]{\linewidth}\raggedright
\textbf{Pages estimées}
\end{minipage} \\
\midrule()
\endhead
\begin{minipage}[t]{\linewidth}\raggedright
\begin{quote}
Page de garde\\
Table des matières
\end{quote}\strut
\end{minipage} & \begin{minipage}[t]{\linewidth}\raggedright
\begin{quote}
Titre, auteur, encadrant, filière, année Générée automatiquement
\end{quote}
\end{minipage} & \begin{minipage}[t]{\linewidth}\raggedright
\begin{quote}
1\\
1
\end{quote}\strut
\end{minipage} \\
\begin{minipage}[t]{\linewidth}\raggedright
\begin{quote}
Introduction\\
Cadre théorique
\end{quote}\strut
\end{minipage} & \begin{minipage}[t]{\linewidth}\raggedright
\begin{quote}
Contexte, problématique, objectifs, plan Fondements mathématiques de
chaque architecture
\end{quote}
\end{minipage} & \begin{minipage}[t]{\linewidth}\raggedright
\begin{quote}
1--2\\
3--5
\end{quote}\strut
\end{minipage} \\
\begin{minipage}[t]{\linewidth}\raggedright
\begin{quote}
Partie expérimentale

Discussion transversale
\end{quote}
\end{minipage} & \begin{minipage}[t]{\linewidth}\raggedright
\begin{quote}
Dataset, prétraitement, architecture, résultats, analyse

Comparaison architectures, adéquation données-modèle
\end{quote}
\end{minipage} & 8--15 (par partie)

2--3 \\
\begin{minipage}[t]{\linewidth}\raggedright
\begin{quote}
Impact sociétal / Éthique Conclusion
\end{quote}
\end{minipage} & \begin{minipage}[t]{\linewidth}\raggedright
\begin{quote}
Implications, biais, limites, contexte local Synthèse, résultats clés,
perspectives
\end{quote}
\end{minipage} & \begin{minipage}[t]{\linewidth}\raggedright
\begin{quote}
1--2\\
1
\end{quote}\strut
\end{minipage} \\
\begin{minipage}[t]{\linewidth}\raggedright
\begin{quote}
Références

Annexes
\end{quote}
\end{minipage} & \begin{minipage}[t]{\linewidth}\raggedright
\begin{quote}
Articles, datasets, frameworks (format IEEE/APA)

Code clé, configurations, résultats complémentaires
\end{quote}
\end{minipage} & 1--2

Variable \\
\bottomrule()
\end{longtable}

\begin{quote}
\textbf{9.2 Critères de qualité rédactionnelle}

• Précision : chaque affirmation technique doit être justifiée (formule,
référence,

résultat expérimental).

• Clarté : éviter les phrases vagues ("le modèle fonctionne bien").
Quantifier

systématiquement.

• Esprit critique : ne pas se limiter à rapporter les résultats.
Analyser les échecs,

discuter les limites, proposer des améliorations.

• Cohérence : le fil conducteur (problématique) doit être visible tout
au long du rapport.

• Références : citer les articles fondateurs des architectures utilisées
(LeCun,

Hochreiter, He, etc.).

\textbf{9.3 Ce qu'il faut éviter}

• Copier-coller du code sans explication. Le code doit être commenté et
justifié dans le

texte.

• Présenter des résultats sans analyse. Un tableau de métriques seul ne
constitue pas

une analyse.

• Ignorer les mauvais résultats. Un modèle qui échoue est aussi
instructif qu'un

modèle performant, à condition d'analyser pourquoi.
\end{quote}

\begin{longtable}[]{@{}
  >{\raggedright\arraybackslash}p{(\columnwidth - 2\tabcolsep) * \real{0.5000}}
  >{\raggedright\arraybackslash}p{(\columnwidth - 2\tabcolsep) * \real{0.5000}}@{}}
\toprule()
\begin{minipage}[b]{\linewidth}\raggedright
•
\end{minipage} & \begin{minipage}[b]{\linewidth}\raggedright
\begin{quote}
Utiliser l'accuracy comme unique métrique sur des classes
déséquilibrées.
\end{quote}
\end{minipage} \\
\midrule()
\endhead
\bottomrule()
\end{longtable}

Page 20

\begin{quote}
\emph{Guide de Référence --- Entraînement de Modèles Deep Learning}
\end{quote}

\begin{longtable}[]{@{}
  >{\raggedright\arraybackslash}p{(\columnwidth - 2\tabcolsep) * \real{0.5000}}
  >{\raggedright\arraybackslash}p{(\columnwidth - 2\tabcolsep) * \real{0.5000}}@{}}
\toprule()
\begin{minipage}[b]{\linewidth}\raggedright
•
\end{minipage} & \begin{minipage}[b]{\linewidth}\raggedright
\begin{quote}
Omettre les détails de reproductibilité : seeds, versions de librairies,
\end{quote}
\end{minipage} \\
\midrule()
\endhead
\bottomrule()
\end{longtable}

\begin{quote}
hyperparammètres.
\end{quote}

Page 21

\begin{quote}
\emph{Guide de Référence --- Entraînement de Modèles Deep Learning}

\textbf{10. Références bibliographiques recommandées}

\textbf{10.1 Articles fondateurs}

{[}1{]} Rosenblatt, F. (1958). The Perceptron: A Probabilistic Model for
Information Storage and Organization in the Brain. Psychological Review,
65(6), 386--408.

{[}2{]} Rumelhart, D. E., Hinton, G. E., \& Williams, R. J. (1986).
Learning representations by back-propagating errors. Nature, 323,
533--536.

{[}3{]} LeCun, Y., Bottou, L., Bengio, Y., \& Haffner, P. (1998).
Gradient-Based Learning Applied to Document Recognition. Proceedings of
the IEEE, 86(11), 2278--2324.

{[}4{]} Hochreiter, S. \& Schmidhuber, J. (1997). Long Short-Term
Memory. Neural Computation, 9(8), 1735--1780.

{[}5{]} Cho, K., et al. (2014). Learning Phrase Representations using
RNN Encoder-Decoder for Statistical Machine Translation. EMNLP 2014.

{[}6{]} He, K., Zhang, X., Ren, S., \& Sun, J. (2016). Deep Residual
Learning for Image Recognition. CVPR 2016.

{[}7{]} Kingma, D. P. \& Ba, J. (2015). Adam: A Method for Stochastic
Optimization. ICLR 2015.

{[}8{]} Srivastava, N., et al. (2014). Dropout: A Simple Way to Prevent
Neural Networks from Overfitting. JMLR, 15, 1929--1958.

{[}9{]} Ioffe, S. \& Szegedy, C. (2015). Batch Normalization:
Accelerating Deep Network Training by Reducing Internal Covariate Shift.
ICML 2015.

{[}10{]} Selvaraju, R. R., et al. (2017). Grad-CAM: Visual Explanations
from Deep Networks via Gradient-based Localization. ICCV 2017.

{[}11{]} Lundberg, S. \& Lee, S. (2017). A Unified Approach to
Interpreting Model Predictions. NeurIPS 2017.

\textbf{10.2 Ressources pédagogiques}

{[}12{]} Goodfellow, I., Bengio, Y., \& Courville, A. (2016). Deep
Learning. MIT Press.

{[}13{]} PyTorch Documentation officielle :
https://pytorch.org/docs/stable/

{[}14{]} PyTorch Tutorials : https://pytorch.org/tutorials/

{[}15{]} Scikit-learn : https://scikit-learn.org/stable/
\end{quote}

Page 22

\end{document}
